\section{Introduction}

Medical licensing examinations such as the \gls{usmle} represent one of the most
demanding knowledge challenges in professional education. Candidates must synthesize
clinical reasoning across dozens of specialties under timed, high-stakes conditions.
Traditional preparation relies on self-directed study, peer quizzing, and expensive
tutoring --- resources not equally accessible to all candidates.

This paper presents \gls{emma}, a conversational medical study agent built to close
that gap. In explanation mode, students pose clinical questions in natural language
and receive responses grounded in passages from eighteen standard medical textbooks.
In quiz mode, \gls{emma} presents authentic \gls{usmle}-style questions, evaluates
answers, and tracks per-specialty performance over time. A collaborative filtering
recommender then steers students toward their weakest areas.

\gls{emma} is built on a five-component \gls{ml} pipeline: a SpaCy \gls{ner} model
for clinical entity extraction, a \gls{faiss} vectorstore over 36,723 textbook chunks,
a \gls{tfidf} + \gls{lsvc} specialty classifier, a BERTopic clustering model for
fine-grained topic routing, and a Surprise-based \gls{cf} recommender. All components
feed a locally-run Qwen3-4B-Thinking-2507 \gls{llm}.

The central empirical question is whether textbook-grounded \gls{rag} improves a
small local \gls{llm}'s accuracy on MedQA compared to a no-retrieval baseline.
A six-run ablation grid shows the answer is condition-dependent: with the right
embedding and \gls{llm} pairing, \gls{rag} delivers an $+11$ \gls{pp}
accuracy gain; without those conditions, it can hurt.

\begin{figure}[ht]
  \centering
  \includegraphics[width=\columnwidth]{pipeline}
  \Description{The EMMA pipeline: NER extracts clinical entities from the query,
    FAISS retrieves relevant textbook chunks, the classifier routes by specialty,
    BERTopic refines by topic cluster, and the LLM generates the final response.}
  \caption{The EMMA system architecture.}
  \label{fig:pipeline}
\end{figure}

The remainder of this paper is organized as follows.
Section~\ref{sec:problem} formalizes the problem.
Sections~\ref{sec:data}--\ref{sec:crs} describe each pipeline component.
Section~\ref{sec:ablation} presents the RAG ablation study.
Section~\ref{sec:evaluation} evaluates results.
Section~\ref{sec:error} analyzes failure modes.
Section~\ref{sec:conclusion} concludes.